%% ABSTRACT (150-250 words, one paragraph, no citations, no \ref).
Black-box REST API testers build their requests from the API's OpenAPI
specification. If the requests do not pass the API's first checks, the tester
spends its budget on rejections and never reaches the code behind them. The
types in the specification are not enough. On the pet-clinic API, a generator
that followed only the types sent \NMotivBaselineRequests requests to add a
pet; all of them failed with a server error while parsing a random date, and
none created a pet. We present \tool, the input component of a testing tool
for the REST League 2027 competition. It parses the specification and
resolves every reference. It fills each request with values taken from
examples, formats, regular-expression patterns and parameter names. It reuses
real ids from earlier responses, and it orders operations so that resources
are created before they are used. On \NSpecs real specifications with
\NOpsTotal operations, \tool left no reference unresolved, and \NConfTool of
its generated values followed their schema, against \NConfBaseline for a
type-only generator. On two competition APIs run live with the same request
budget, it raised the operations that returned a success from
\NPetclinicBaselineCov to \NPetclinicFullCov and from
\NGestaohospitalBaselineCov to \NGestaohospitalFullCov. Measuring each part
on its own showed that schema-aware values and id reuse gave this gain, while
ordering operations by their dependencies changed the result by less than one
operation on either API.
